\documentclass[12pt,a4paper]{article}
\usepackage{cmap} % Makes the PDF copiable. See http://tex.stackexchange.com/a/64198/25761
\usepackage[T1]{fontenc}
\usepackage[brazil]{babel}
\usepackage[utf8]{inputenc}
\usepackage{amsmath}
\usepackage{amsfonts}
\usepackage{amssymb}
\usepackage{amsthm}
\usepackage{textcomp} % \degree
\usepackage{gensymb} % \degree
\usepackage[usenames,svgnames,dvipsnames]{xcolor}
\usepackage{hyperref}
\usepackage{graphicx}
\usepackage[margin=2cm]{geometry}

\hypersetup{
    colorlinks = true,
    allcolors = {blue}
}

% TODO: Consider using exsheets
% http://linorg.usp.br/CTAN/macros/latex/contrib/exsheets/exsheets_en.pdf
%
% http://ctan.org/tex-archive/macros/latex/contrib/exercise/
% Options: answerdelayed,lastexercise,noanswer
\usepackage[answerdelayed,lastexercise]{exercise}

\addto\captionsbrazil{%
\def\listexercisename{Lista de exerc\'icios}%
\def\ExerciseName{Exerc\'icio}%
\def\AnswerName{Solu\c{c}\~ao do exerc\'icio}%
\def\ExerciseListName{Ex.}%
\def\AnswerListName{Solu\c{c}\~ao}%
\def\ExePartName{Parte}%
\def\ArticleOf{de\ }%
}

\renewcommand{\ExerciseHeaderTitle}{(\ExerciseTitle)\ }
\renewcommand{\ExerciseListHeader}{%\ExerciseHeaderDifficulty%
\textbf{%\ExerciseListName\
\ExerciseHeaderNB.\ %
%\ --- \
\ExerciseHeaderTitle}%
%\ExerciseHeaderOrigin
\ignorespaces}
\renewcommand{\AnswerListHeader}{\textbf{\ExerciseHeaderNB.\ (\AnswerListName)\ }}

\newcommand{\fixme}{{\color{red}(...)}}
\newcommand*\sen{\operatorname{sen}}
\newcommand*\R{\mathbb{R}}

\renewcommand{\theenumi}{\alph{enumi}}
\renewcommand\labelenumi{(\theenumi) }

\newcommand*\tipo{Prova II}
\newcommand*\turma{TADS121-01U}
\newcommand*\disciplina{CDI0001}
\newcommand*\eu{Helder G. G. de Lima}
\newcommand*\data{23/10/2017}

\author{\eu}
\title{\tipo - \disciplina}
\date{\data}

\begin{document}
\thispagestyle{empty}
\newgeometry{margin=2cm,bottom=0.5cm}
\begin{center}
\includegraphics[width=9.0cm]{marca} \\
\textbf{\tipo\ (\disciplina / \turma)} \\
Prof. \eu\footnote{
Este é um material de acesso livre distribuído sob os termos da licença \href{https://creativecommons.org/licenses/by-sa/4.0/deed.pt_BR}{Creative Commons BY-SA 4.0}}
\end{center}

\noindent Nome do(a) aluno(a): \underline{\hspace{9,7cm}} Data: \underline{\data}

\begin{center}\fbox{
\begin{minipage}{14cm}

{\footnotesize
\begin{itemize}
\renewcommand{\theenumi}{\Roman{enumi}}
\item Identifique-se em todas as folhas.
\item Mantenha o celular e os demais equipamentos eletrônicos desligados durante a prova.
\item Justifique cada resposta com cálculos ou argumentos baseados na teoria estudada.
\item Escolha e responda apenas 5 das questões, totalizando 10,0 pontos.
\end{itemize}
}

\end{minipage}
}
\end{center}

\begin{ExerciseList}
\Exercise[title={2,0}]
Utilize a definição $\displaystyle f^\prime(x_0) = \lim_{x \to x_0} \dfrac{f(x) - f(x_0)}{x-x_0}$ para calcular a derivada de uma função $f$ de sua escolha, no ponto $x_0 = 2$.
\Answer \fixme

\Exercise[title={2,0}] Use as regras de derivação para obter uma fórmula (devidamente simplificada) para as derivadas das seguintes funções:
\begin{enumerate}
\item $s(x) = 3x - \left(\dfrac{1}{x}\right)^{14}+23\sen(4x)$
\item $q(x) = \dfrac{e^{2x}+e^{4x}+6}{3e^x}$
\end{enumerate}
\Answer
\begin{enumerate}
\item Como $\left(\dfrac{1}{x}\right)^{14} = \dfrac{1}{x^{14}} = x^{-14}$, tem-se
\[
s^\prime(x)
= (3x)^\prime - (x^{-14})^\prime + (23\sen(4x))^\prime
= 3 + 14x^{-15}+23 \cos(4x) \cdot 4
= 3 + \dfrac{14}{x^{15}} + 92\cos(4x).
\]
\item Como $q(x)
= \dfrac{e^{2x}+e^{4x}+6}{3e^x}
= \dfrac{e^{2x}}{3e^x} + \dfrac{e^{4x}}{3e^x} +\dfrac{6}{3e^x}
= \dfrac{1}{3}(e^{x}+e^{3x}+6e^{-x})
$, tem-se
\[
q^\prime(x)
= \dfrac{1}{3}\left[(e^{x})^\prime+(e^{3x})^\prime+6(e^{-x})^\prime\right].
= \dfrac{1}{3}\left[e^{x} + 3e^{3x} - 6e^{-x} \right]
= \dfrac{e^x}{3} + e^{3x} - 2e^{-x}.
\]
\end{enumerate}


\Exercise[title={2,0}] Utilizando as regras de derivação estudadas, obtenha a equação da reta $y = ax + b$ que tangencia o gráfico de $h(x) = \ln(\sqrt{x+2})$ no ponto $P = (x_0, y_0)$ cuja abscissa é $x_0 = -1$.
\Answer \fixme

\Exercise[title={2,0}] Considere a função $g: \R \to \R$ definida por $g(x) = x e^x$. Mostre que as derivadas de ordem superior de $g$ satisfazem o seguinte: $g^{\prime}(-1) = g^{\prime\prime}(-2) = g^{\prime\prime\prime}(-3) = 0$.
\Answer \fixme

\Exercise[title={2,0}] Seja $y = f(x)$ uma função derivável dada \textit{implicitamente} pela equação $xy + \ln(y) = 0$.
Calcule $y'$ e determine se há algum ponto do gráfico de $f$ em que a reta tangente seja horizontal.
\Answer \fixme
\Exercise[title={2,0}] Uma esfera aumenta de modo que seu raio cresce à razão de $2 cm/s$. Qual é a taxa de variação do volume no instante em que o raio é de $10cm$?
\Answer \fixme

\end{ExerciseList}

\begin{center}
BOA PROVA!
\end{center}

%\newpage
%\restoregeometry
%\section*{Respostas}
%\shipoutAnswer
\end{document}
